\documentclass[12pt]{article}
\usepackage[utf8]{inputenc}
\usepackage[T1]{fontenc}
\usepackage[margin=1in]{geometry}
\usepackage{setspace}
\usepackage{xcolor}
\usepackage[colorlinks=true, urlcolor=blue!60!black, linkcolor=blue!60!black]{hyperref}
\setlength{\parindent}{0em}
\setlength{\parskip}{0.6em}
\renewcommand{\baselinestretch}{1.0}
\begin{document}
\begin{center}
{\LARGE \textit{INS v. Chadha}, 462 U.S. 919 (1983)\par}
\vspace{0.5em}
{\large Case Brief --- Two-Page Edition\par}
\end{center}

\section{Memory Jogger}
The one-House legislative veto is unconstitutional: once Congress delegates power, it can alter rights outside the Legislative Branch only through bicameralism and presentment (\href{https://supreme.justia.com/cases/federal/us/462/919/}{462 U.S. 919}).

\section{Detailed Case Facts}
Jagdish Rai Chadha overstayed a student visa, and an Immigration Judge suspended his deportation for ``extreme hardship'' under \S~244 of the Immigration and Nationality Act. Section 244(c)(2) let either House veto such a suspension. In 1975 the House, without debate or a recorded vote, passed a resolution disapproving suspension for Chadha and five others; it never went to the Senate or the President.

\section{Procedural History}
Chadha was ordered deported; the Ninth Circuit, with the House and Senate participating, held the veto unconstitutional. The Supreme Court affirmed.

\section{Judicial Votes}
Affirmed, 7--2: Burger wrote for Brennan, Marshall, Blackmun, Stevens, and O'Connor; Powell concurred in the judgment; White and Rehnquist dissented.

\section{Holding}
The case was justiciable and the veto severable. The one-House veto was legislative action that failed the bicameralism and presentment requirements of Article I, \S\S~1 and 7.

\section{Analysis of Opinions}
Burger held that efficiency cannot save a procedure contrary to constitutional text. Action is legislative when it alters the legal rights and duties of persons outside the Legislative Branch, and the Constitution's express unicameral exceptions foreclose others; Congress ``must abide by its delegation'' until it legislates. Powell would have held only that the House usurped a judicial function. White defended the veto as an ``indispensable, political invention'' for the administrative state, sparing Congress a ``Hobson's choice''; Rehnquist found the veto inseverable.

\section{Examples: Future Applications}
\textbf{Same-Side:} a two-House concurrent resolution disapproving agency rules is invalid, as in \textit{Process Gas}. \textbf{Opposite-Side:} a report-and-wait scheme allowing disapproval only by a presented joint resolution, like the Congressional Review Act, is valid. \textbf{Fence-Sitter:} a requirement that both Houses and the President approve major rules before they take effect.

\section{Critique}
Footnote 16 treats agency rulemaking as executive and congressional review as legislative, though both alter rights outside Congress. Severing the veto arguably gave the Executive broader unchecked power than Congress ever delegated.

\section{Key Quotations}
\begin{quote}
``It emerges clearly that the prescription for legislative action in Art. I, \S\S~1, 7, represents the Framers' decision that the legislative power of the Federal Government be exercised in accord with a single, finely wrought and exhaustively considered, procedure.''
\end{quote}
\noindent --- Chief Justice Burger, Opinion of the Court, 462 U.S. at 951.

\end{document}
